\documentclass[12pt,a4paper]{article}
\usepackage{enumitem}
\usepackage[utf8]{inputenc}
\usepackage[spanish]{babel}
\usepackage{csquotes}
\usepackage{setspace}
\usepackage{geometry}
\usepackage{enumitem}
\setlist[enumerate,1]{label=\arabic*.}
\setlist[enumerate,2]{label*=\arabic*.}
\setlist[enumerate,3]{label*=\arabic*.}

\usepackage{hyperref}
\geometry{margin=2.5cm}

\setstretch{1.2}

\title{CONSIDERACIONES PARA LA ELABORACIÓN DE LA TESIS \\ Estilo Chicago}
\author{Staling Cordero-Brito \\ (Con ejemplos ilustrativos de las Bibliotecas UAI -- AD)}
\date{}

\begin{document}
\maketitle

\section{Reglas de estilo en el trabajo escrito (Maestría de Historia y Geografía para Educadores)}

\subsection*{Tipografía}
\begin{enumerate}[label=\alph*)]
    \item La tipografía de la letra de la tesis debe ser Times New Roman, tamaño 12, con interlineado de un espacio y medio (1.5), y dos espacios (2.0) entre párrafos y entre títulos y subtítulos.
    \item Todos los nombres de los capítulos van en letra 12, en negrita y en mayúsculas sostenidas.
    \item Todos los títulos van en letra 12, en negrita y con letra mayúscula inicial, es decir, con la primera letra en mayúscula, salvo en el caso de los nombres propios, que siempre inician con letra mayúscula.
    \item El cuerpo del documento se trabajará en color negro, a excepción de los gráficos y tablas.
    \item Las negritas solo se utilizarán para los nombres de capítulos, los títulos y los subtítulos.
    \item Los gráficos y tablas se colocarán con sus títulos y fuentes correspondientes.
    \item Las palabras que deben llevar acento, aun en mayúsculas, deberán estar acentuadas.
    \item Todo el cuerpo del trabajo va justificado a ambos lados. Sin embargo, las citas de más de 40 palabras se escribirán con sangría de 10 espacios y con interlineado sencillo.
\end{enumerate}
\newpage
---
\section{Ejemplo de estructura de la tesis}

\begin{center}
\textbf{Índice}
\end{center}

% --- Parte preliminar (árabe simple) ---

\noindent I. Dedicatoria\\
II. Agradecimientos\\
III. Resumen (palabras clave)\\
IV. Introducción\\

% --- Capítulo I ---
\noindent\textbf{Capítulo 1: Planteamiento del problema}
\begin{enumerate}[resume] % continúa desde 3 → ahora será 4, pero lo ajustamos
  \setcounter{enumi}{0}   % reinicia para que el capítulo empiece en 1

  \item Formulación del problema
  \item Delimitación
  \item Justificación
  \item Alncance
  \item Preguntas de investigación
  \item Hipótesis (si aplica)
  \item Objetivos
  \begin{enumerate}
      \item General
      \item Específicos
  \end{enumerate}
  \item Operacionalización de las variables
\end{enumerate}

% --- Capítulo II ---
\noindent\textbf{Capítulo 2: Marco teórico de la investigación}
\begin{enumerate}[resume]
  \item Antecedentes
   \begin{enumerate}
      \item Antecedentes históricos (Evolución del tema a lo largo del tiempo)
      \item Antecedentes internacionales (Investigaciones internacionales sobre el tema por diversos autores)
      \item Antecedentes nacionales (Investigaciones nacionales sobre el tema por diversos autores)
  \end{enumerate}
   \item Revisión de la literatura (Perspectiva teórica)
\end{enumerate}

% --- Capítulo III ---
\noindent\textbf{Capítulo 3: Diseño metodológico de la investigación}
\begin{enumerate}[resume]
  \item Tipo de investigación
  \item Enfoque de la investigación
  \item Métodos de estudios
  \item Técnicas de recolección de datos
  \item Instrumentos de recolección de datos
    \item Validez y confiabilidad de los instrumentos
    \item Universo, Población y muestra

% --- Capítulo IV ---
\noindent\textbf{Capítulo 4: Presentación de los resultados}
\begin{enumerate}[resume]
 Discusión \\
  Conclusiones\\
  Recomendaciones\\
  Referencias Bibliográficas\\
  Anexos
\end{enumerate}

\vspace{0.5cm}
\textit{Nota: Esta estructura es solo un ejemplo orientativo. Puede variar el tipo de trabajo de cada maestrante.}

---

\section{Introducción al estilo Chigago}
El formato Chicago se caracteriza por integrar dos métodos principales de citación: el uso de notas al pie que posteriormente se sistematizan en una sección final de referencias o bibliografía, y la inclusión de citas breves dentro del texto.\\

Este sistema busca facilitar el acceso a la información bibliográfica para autores, editores y otros profesionales vinculados a la producción académica, promoviendo una experiencia clara y eficiente en la gestión de fuentes.\\

La citación en Chicago se fundamenta en dos modalidades: por un lado, las notas y la bibliografía, y por otro, las citas breves en el texto que incorporan el apellido del autor y el año de publicación, diferenciadas por el uso de paréntesis.\\

El método de notas y bibliografía es especialmente frecuente en disciplinas humanísticas como literatura, historia, artes y filosofía, donde las fuentes se organizan mediante numeración y pueden ubicarse al pie de página o al final de cada sección, permitiendo agregar comentarios adicionales según el contexto de la cita.\\

En contraste, el sistema autor-fecha es predominante en las ciencias sociales, presentando la información de manera concisa dentro del texto e indicando el autor y la fecha de la obra citada. Cada referencia en el texto se vincula con una entrada detallada en la lista de referencias o bibliografía.\\

Ambos enfoques comparten criterios de presentación similares, por lo que se recomienda consultar publicaciones especializadas del área o seguir las directrices editoriales correspondientes.\\

La última edición del Manual de Estilo Chicago (17ª) incorpora directrices actualizadas para citar fuentes digitales, considerando aspectos como la inmediatez informativa, el uso de dispositivos electrónicos, la gestión de metadatos y el trabajo colaborativo en entornos virtuales.\\


\subsection{Objetivos de aprendizaje}
Al finalizar la lectura de esta Guía, se espera que el lector sea capaz de identificar los aspectos centrales del estilo de citación Chicago, al mismo tiempo que pueda establecer sus necesidades de información y ubicar las respuestas dentro de la estructura de contenidos presentada a continuación.

\subsection{Citas y fuentes}
Las citas y entradas de fuentes son conceptos que suelen confundirse a la hora de aplicar un estilo de citación determinado cuando se elabora un texto. 

En este sentido, muchas veces se utilizan los términos \enquote{cita}, \enquote{bibliografía} y \enquote{referencia} como símiles, cuando en realidad cada uno corresponde a un elemento de la estructura de citación. 

\subsection*{Diferencias básicas}
\begin{itemize}
    \item \textbf{Cita:} Mención al texto del cual se extrajo una idea o frase ajena, permitiendo al lector distinguir entre el contenido propio y aquel utilizado para apoyar la línea argumentativa presentada.
    \item \textbf{Referencia:} Incluye los elementos que permiten la identificación de la fuente originaria.
    \item \textbf{Bibliografía:} Relación de textos utilizados como fuentes documentales o apoyo, no necesariamente citados en el texto.
\end{itemize}

\subsection{Sistema Autor-Fecha}
Este estilo se usa principalmente en ciencias y ciencias sociales. Está compuesto por el apellido del autor, el año de publicación y, en caso de ser relevante, el número de página.\\

El formato de citación está compuesto por el nombre del/los autor/es, el año de publicación y, en caso de ser relevante, el número de página. Este conjunto de datos corresponde a una cita, para lo cual cada una debe contar con una entrada en el listado de referencias dispuesto
al final del manusrito, documento o artículo. Es importante destacar que cada entrada debe contener todos los detalles de la fuente.

\subsection*{Ejemplo de cita en texto}
\begin{quote}
(Koontz 2010, 48)
\end{quote}

\subsection*{Ejemplo de entrada en lista de referencias}
\begin{quote}
Koontz, Harold. 2010. \textit{Administración: una perspectiva global}. Madrid: McGraw Hill.
\end{quote}

\subsubsection*{Ejemplos de citas en el sistema autor-fecha}

Como se muestra en el ejemplo, este tipo de cita se ubica en el texto utilizando paréntesis, para lo cual se agrega el apellido del autor y la fecha, sin puntuación entre ambos datos:

\begin{quote}
(Cordero-Brito 2024)
\end{quote}

Si la cita refiere a una parte específica del texto utilizado como fuente, o parafraseo, se recomienda agregar el número de página ---o rango de páginas separados por un guion--- posterior a la fecha seguida por una coma:

\begin{quote}
(Cordero-Brito 2024, 122--215)
\end{quote}

En cuanto a su ubicación, las citas generalmente se colocan al final de una idea importante, oración o cita textual, antes de la puntuación.  
En los casos en que se requiera citar más de una fuente en una misma oración, la cita debe incluir ambas fuentes dentro de un paréntesis, pero separadas por un punto y coma:

\begin{quote}
Ante ello, la transmisión de flujos informativos ofrece a quien la recibe un espacio de libertad y socialización universal (Rivera 2002). Por tanto, la pregunta radicará no solo en cuestionarse cómo aprendieron los alumnos lo que se les enseña, sino qué hacen con lo que se les enseñó (Acosta 1998; Roa 2001).
\end{quote}

Si el nombre del autor citado está incluido en la redacción del texto, la citación debe incluir sólo la fecha luego de mencionar el apellido.  
Mientras que para las citaciones textuales se debe incluir también el número de página entre paréntesis:

\begin{quote}
Para Cordero-Brito (2023) es importante conocer el tipo de jugador que tenemos en el aula para determinar los motivadores que voy a utilizar en el aula ``Se clasifican K, A, S E'' (1O), por ejemplo, la taxonomaía de Bartle.
\end{quote}



\subsection{Sistema Notas y Bibliografía}


Para utilizar el estilo de notas y bibliografía de Chicago se debe utilizar la opción de número a pie de página (disponibles en los editores de texto) y colocarlo al final de la oración, o de las comillas que encierran a la cita textual, y antes de la puntuación:

\begin{quote}
Y también, que existen algunos atributos que suponen antagonismos en la actualidad entre los roles de mujer y científica, pues ``la mujer científica de hoy tiene características que rompen su rol tradicional, es ambiciosa, competitiva y descuida su familia''\textsuperscript{1}.
\end{quote}

\rule{\linewidth}{0.4pt}

\noindent \textsuperscript{[Gies, \textit{creencias sobre ciencia}, 18.] o [Erica Gies,  .a mujuer en ciencia: retos y perspectivas (University of Chicago Press, 2022), 27.]
\vspace{1em}}

Las ``notas a pie'' se visualizan al final de cada página, mientras que las ``notas finales'' al final del texto. Si se desea seguir con el estilo de notas de Chicago, se debe escoger entre el uso de notas al pie o notas finales, y usarlo de manera consistente, sin mezclar entre ambos tipos de notas.


\subsection*{Ejemplo de cita en nota completa}
\begin{quote}
1. Xavier Zubiri, \textit{Inteligencia Artificial} (España: Tecnos, 2004), 58--63.
\end{quote}

\subsection*{Ejemplo de nota corta}
\begin{quote}
Zubiri, \textit{Inteligencia Artificial}, 58--63.
\end{quote}

\section{Listado de referencias y Bibliografía}
\begin{itemize}
    \item El listado de referencias se utiliza sólo en el sistema Autor-fecha.
    \item La bibliografía se utiliza sólo en el sistema de notas.
\end{itemize}

\subsection*{Ejemplo de referencia (autor-fecha)}
\begin{quote}
Bernstein, Bernard. 1994. \textit{La estructura del discurso pedagógico}. Madrid: Morata S.A.
\end{quote}

\subsection*{Ejemplo de bibliografía (notas)}
\begin{quote}
Zubiri, Xavier. \textit{Inteligencia sentiente}. España: Tecnos, 2004.
\end{quote}

\subsection{Tipos de referencias}
A continuación, se muestran ejemplos de referencias según el tipo de material.

\subsubsection{Libro}
\textbf{Autor-fecha:}  
(Zubiri 2004, 58--63)  

\textbf{Referencia:}  
Zubiri, Xavier. 2004. \textit{Inteligencia sentiente}. España: Tecnos.  

\textbf{Nota completa:}  
Xavier Zubiri, \textit{Inteligencia sentiente} (España: Tecnos, 2004), 58--63.  

\textbf{Bibliografía:}  
Zubiri, Xavier. \textit{Inteligencia sentiente}. España: Tecnos, 2004.  

\subsubsection{Artículo de revista}
\textbf{Autor-fecha:}  
(Gaete 2011, 12)  

\textbf{Referencia:}  
Gaete, Marcela. 2011. ``Articulación del sistema de educación superior en Chile: posibilidades, tensiones y desafíos''. \textit{Calidad en la Educación}, 51, no. 35 (primavera): 22--51. DOI 10.31619/caledu.n35.96.  

\textbf{Nota completa:}  
Marcela Gaete, ``Articulación del sistema de educación superior en Chile: posibilidades, tensiones y desafíos'', \textit{Calidad en la Educación}, 51, no. 35 (abril 2011): 22--51. DOI 10.31619/caledu.n35.96.  

\textbf{Bibliografía:}
Gaete, Marcela. 2011. “Articulación del sistema de educación superior en Chile: 
posibilidades, tensiones y desafíos”. \textit{Calidad en la Educación}, 51, no. 35 
(primavera): 22--51. DOI: 10.31619/caledu.n35.96.

\subsubsection{Periódico}
\textbf{Autor-fecha:}  
(Kuchler 2021)  

\textbf{Referencia:}  
Kuchler, Hannah. 2021. ``Vaccitech warns of blood clot risk in IPO filing''. \textit{Financial Times}, 09 abril, 2021. \url{https://www.ft.com/content/06860275-4cbe-4cf3-9251-5e47e87ccf0d}  

\textbf{Notas y bibliografía}

\noindent\textbf{Nota completa:}  
Hannah Kuchler, “Vaccitech warns of blood clot risk in IPO filing”, \textit{Financial Times}, 09 abril, 2021, \url{https://www-ft-com.uai.idm.oclc.org/content/06860275-4cbe-4cf3-9251-5e47e87ccf0d}.

\noindent\textbf{Nota corta:}  
Kuchler, “Vaccitech warns of blood clot risk in IPO filing”.

\noindent\textbf{Entrada de bibliografía:}  
Kuchler, Hannah. “Vaccitech warns of blood clot risk in IPO filing”. \textit{Financial Times}, 09 abril, 2021.  
\url{https://www-ft-com.uai.idm.oclc.org/content/06860275-4cbe-4cf3-9251-5e47e87ccf0d}


\subsubsection{Sitio web} 
\textbf{Autor-fecha:}  
(Universidad Adolfo Ibáñez 2021)  

\textbf{Referencia:}  
Universidad Adolfo Ibáñez. 2021. ``Admisión UAI''. Recuperado el 04 de abril, 2021. \url{https://www.uai.cl/admision}  

\textbf{Notas y bibliografía}

\noindent\textbf{Nota completa:}  
“Admisión UAI”, \textit{Universidad Adolfo Ibáñez}, recuperado el 04 de abril, 2021, \url{https://www.uai.cl/admision}.

\noindent\textbf{Nota corta:}  
“Admisión UAI”.

\noindent\textbf{Entrada de bibliografía:}  
Universidad Adolfo Ibáñez. ``Admisión UAI''. Recuperado el 04 de abril, 2021.  
\url{https://www.uai.cl/admision}\\


\textbf{Staling Cordero-Brito} \\
Maestría en Historia y Geografía para Educadores \\
Universidad Autónoma de Santo Domingo (UASD) \\
Correo:scordero48@uasd.edu.do

\end{document}
